\chapter{Ce qui existe déjà}
\label{ch:etat-existant}

Ce chapitre fait le tour des outils et méthodes qui traitent
aujourd'hui la conformité d'un SI. L'objectif n'est pas d'être
exhaustif --- les références bibliographiques le sont --- mais de
situer où se place HOMO-CI par rapport à ce qui existe.

\section{Le cadre ANSSI : un dossier daté}

L'ANSSI publie un \emph{Guide d'homologation en neuf étapes
simples}~\cite{anssi-homologation-2014}. Les neuf étapes vont du
cadrage (étape~1) à la décision d'homologation (étape~8), avec
une étape~9 souvent oubliée~: \emph{maintenir l'homologation dans
le temps}.

La méthode pour analyser le risque, c'est \textbf{EBIOS Risk
Manager}~\cite{anssi-ebios-rm-2018}. Cinq ateliers~: cadrage,
sources de risque, scénarios stratégiques, scénarios
opérationnels, traitement du risque. EBIOS RM fournit la matière
de plusieurs sections du dossier (analyse de risque, mesures de
sécurité).

Le point commun de ces deux outils, c'est qu'ils produisent un
livrable \emph{ponctuel}. Un dossier daté, une analyse de risque
datée. Rien dans le cadre actuel ne dit comment ces livrables
sont mis à jour quand le SI évolue, à part une re-conduction
complète tous les 3~ans.

\section{L'audit annuel : ce qui se pratique}

Dans la plupart des organisations publiques, le suivi de
l'homologation se résume à~:

\begin{itemize}
  \item un audit \textbf{annuel} ou tri-annuel, externe ou
    interne, qui regarde un échantillon de règles et de
    configurations~;
  \item une re-conduction de l'homologation tous les 3~ans, qui
    mobilise typiquement entre 40 et 120 jours-homme selon la
    complexité du SI.
\end{itemize}

Le coût est élevé, et entre deux audits, on est aveugle. Si une
règle NSG est désactivée par erreur en juin, on s'en rendra
compte (peut-être) à l'audit de l'année suivante.

\section{Le DevSecOps et la \emph{continuous compliance}}

Côté privé et cloud, depuis 2017--2018, des outils proposent une
collecte continue de preuves~: \textbf{Drata}, \textbf{Vanta},
\textbf{Tugboat Logic} pour les certifications SOC 2 / ISO 27001
américaines. Le principe est le même que ce qu'on fait ici~:
brancher des collecteurs sur les API du SI, regarder en
permanence si les contrôles sont respectés.

Deux limites pour notre cas~:

\begin{enumerate}
  \item Ces outils ciblent SOC 2 et ISO 27001, pas le cadre ANSSI.
    Le mapping vers l'homologation ANSSI doit être fait à la main.
  \item Ce sont des produits commerciaux fermés. On ne peut ni
    auditer leur code, ni l'adapter aux spécificités françaises
    (EBIOS RM, AQSSI, dossier au format ANSSI).
\end{enumerate}

\section{Les outils d'inventaire et de posture}

Côté cloud, des outils comme \textbf{Azure Policy},
\textbf{AWS Config}, \textbf{Prowler}, \textbf{ScoutSuite}
collectent l'état des ressources et le comparent à des
référentiels. Ils sont utiles, mais s'arrêtent à l'inventaire~:
ils ne produisent pas le dossier d'homologation, ni l'analyse de
risque associée. Ils sont une brique~; HOMO-CI les utilise (le
watcher NSG s'appuie sur l'API Azure Resource Manager), mais va
plus loin en produisant un livrable réglementaire.

\section{Où se place HOMO-CI}

\begin{table}[h]
\centering
\caption{Positionnement de HOMO-CI par rapport à l'existant.}
\label{tab:positionnement}
\small
\begin{tabularx}{\textwidth}{lXcccc}
\toprule
\textbf{Outil} & \textbf{Cible} & \textbf{Continu} & \textbf{ANSSI} & \textbf{Open source} & \textbf{Dossier complet} \\
\midrule
Audit annuel manuel & ANSSI / ISO & \xmark & \cmark & --- & \cmark \\
Drata, Vanta & SOC 2 / ISO 27001 & \cmark & \xmark & \xmark & partiel \\
Prowler, ScoutSuite & Posture cloud & \cmark & \xmark & \cmark & \xmark \\
Azure Policy & Conformité Azure & \cmark & \xmark & \xmark & \xmark \\
\textbf{HOMO-CI} & \textbf{ANSSI} & \cmark & \cmark & \cmark & \cmark \\
\bottomrule
\end{tabularx}
\end{table}

L'apport ne tient pas dans une innovation algorithmique~: les
briques (Wazuh, API Azure, Trivy, Jinja2) existent et sont
matures. L'apport tient dans \textbf{l'intégration} de ces
briques pour produire un livrable conforme au cadre ANSSI, et
dans \textbf{la démonstration} que ça tient debout sur un SI de
type ENT universitaire.

\begin{takeaway}
Le marché propose soit des outils ponctuels et chers (audit
manuel ANSSI), soit des outils continus mais sur d'autres
référentiels (SOC 2, ISO~27001). Personne ne fait du continu
sur ANSSI en open source. C'est le créneau de HOMO-CI.
\end{takeaway}
